%% ma: L12-B06-0034
%% lop: 12
%% chuong: 2
%% bai: 6
%% dang: 12.6.15
%% loai: DS
%% muc: [TH, NB, TH, TH]
%% dap_an: "SĐĐĐ"
%% nguon: "LaTeX giáo viên gửi 21/10/2026 (M05, Vectơ và hệ trục tọa độ trong không gian, 107 câu), Phần II Câu 8"
%% kien_thuc: "Bài 6 (tích vô hướng, ứng dụng tính công)"
%% trang_thai: cho-duyet
%% ly_do: "Dạng mới đề xuất 12.6.15"
\begin{ex}
Một chiếc tàu kéo một xà lan trên mặt biển bằng một lực kéo $\overrightarrow{F}$ có độ lớn $5000$ N. Sợi dây cáp nối tàu và xà lan tạo với phương ngang (hướng dịch chuyển) một góc $30^\circ$. Biết xà lan di chuyển thẳng đều được một quãng đường dài $4$ km.
\choiceTF
{Hình chiếu của lực kéo $\overrightarrow{F}$ lên phương dịch chuyển có độ lớn là $2500$ N.}
{\True Công do lực kéo $\overrightarrow{F}$ sinh ra có giá trị dương.}
{\True Công sinh ra bởi lực $\overrightarrow{F}$ trong quá trình trên là $10^7\sqrt3$ (J).}
{\True Nếu dây cáp nằm song song với mặt biển (góc $0^\circ$) thì công của lực kéo đạt giá trị lớn nhất (với cùng độ lớn lực kéo và cùng quãng đường).}
\loigiai{a) Độ lớn hình chiếu: $5000\cos30^\circ=2500\sqrt3\approx4330$ (N). Sai.

b) $A=F\cdot d\cdot\cos30^\circ>0$. Đúng.

c) $A=5000\cdot4000\cdot\dfrac{\sqrt3}2=10^7\sqrt3$ (J). Đúng.

d) $A=F\cdot d\cdot\cos\alpha$ lớn nhất khi $\cos\alpha=1$, tức $\alpha=0^\circ$. Đúng.}
\end{ex}
